\documentclass[11pt]{article}
\usepackage[T1]{fontenc}
\usepackage[utf8]{inputenc}
\usepackage{amsmath,amssymb,booktabs,array,longtable}
\usepackage[margin=1in]{geometry}
\usepackage[hidelinks]{hyperref}
\usepackage{url}
\hypersetup{pdftitle={Supplementary material: Reproducibility guide for R(F3,F4)},
  pdfauthor={Hanxin Jiang and Zehui Shao}}
\setlength{\parindent}{0pt}
\setlength{\parskip}{0.45em}
\setlength{\emergencystretch}{2em}
\newcommand{\file}[1]{\path{#1}}
\title{Supplementary Material\\[0.4em]
\large Reproducibility guide for an $8+9$ gluing computation\\
for the Ramsey number $R(F_3,F_4)$}
\author{Hanxin Jiang \quad Zehui Shao\\[0.4em]
\small Institute of Computing Science and Technology,\\
\small Guangzhou University, Guangzhou 510006, China}
\date{}
\begin{document}
\maketitle

\section{Purpose and verification scope}

This supplement explains how to compile the search program, reproduce the
seed catalogues and filtering counts, run the distributed search, and verify
the certificates for the final forty seed pairs. Here $F_t$ is the graph
formed by $t$ triangles with a common vertex and otherwise disjoint vertex
sets. A colouring is admissible if it contains neither a red $F_3$ nor a blue
$F_4$. The mathematical justification of the reductions is given in the
accompanying article.

The bulk-search results and cumulative solver counts reported in the
article were obtained by running the C++ search program on the authors'
computing cluster. The local certificate checks described below were
performed separately after the distributed search. Thus the cluster
statistics and the subsequent independent verification results have
distinct computational provenance.

The archived search has $88\,294\,674$ retained ordered pairs of seeds on
eight and nine vertices. The reported bulk computation excludes
$88\,294\,634$ of these pairs. The remaining forty are excluded separately:
thirty-four by the counting argument in the article, with its finite
premises checked from the seed graphs, and six by independently replayable
finite-domain certificates. The quick verification below establishes these
forty local exclusions. It does not rerun or certify the other
$88\,294\,634$ solver decisions or prove catalogue completeness by itself.

The distribution provides the inputs and source needed for a fresh full
computation. It does not include the complete historical cluster
checkpoints or a checked propositional proof for each bulk exclusion.
The supplied source, the completeness arguments in the article, and a
full rerun therefore have different roles from the local certificate
checker. File checksums establish file identity, not mathematical truth.

\section{Files and conventions}

Extract \file{fan_F3F4_reproducibility.zip} to a writable directory. All
commands below are issued from the extracted directory containing
\file{README.txt}, unless a different working directory is stated.
The shell examples use Linux syntax. Section~\ref{sec:build} gives the
corresponding Windows build procedure.

\begin{center}
\begin{tabular}{@{}p{0.25\linewidth}p{0.68\linewidth}@{}}
\toprule
Location & Contents and purpose \\
\midrule
\file{src/} & C++17 search source, build files, and required third-party
source and licences. \\
\file{data/} & Fixed graph6 seed files, the ordered-pair catalogue, and
the dataset identity. \\
\file{tools/} & Canonical generation, catalogue preparation, and
independent data-validation programs. \\
\file{certificates/} & The forty case specifications and final
counting and finite-domain certificate data. \\
\file{verify40.py} & Entry point for independent verification of the
forty exceptional pairs. \\
\file{combine_results.py} & Checks coverage of a fresh full run and
combines it with the verified local exclusions. \\
\file{docs/} & This supplement in editable LaTeX and PDF form. \\
\bottomrule
\end{tabular}
\end{center}

A seed index and a case index are zero-based. Machine numbers are one-based.
The red graph determines the colouring; all other edges of the complete
graph are blue. A graph6 record encodes an undirected graph, using the
vertex order fixed by the corresponding seed file. The production
dataset identifier is \texttt{e38c0723b6a3be74}; the engine identifier is
\texttt{fan17-cpp-1}. The eight-vertex and nine-vertex sides are ordered:
there is no transposition identification between them.

The published files should be kept unchanged. Put new checkpoints,
timing logs, generated certificates, and validation reports in separate
working directories. These outputs are produced by the reader's run and
are not prerequisites for the archived quick check.

\section{Requirements}

The forty-case checker requires Python 3.9 or later and only its standard
library. The C++ search requires a C++17 compiler and threads. GCC 9 or
later, a recent Clang, or the x64 Microsoft Visual C++ toolchain can be
used. A CMake build requires CMake 3.15 or later. The included direct
build scripts provide an alternative to CMake. MiniSat and the JSON
header are supplied in the archive; no separate SAT-solver installation
is needed.

Full seed and profile regeneration requires Python 3.11 or later and
the pinned Python packages in \file{tools/requirements.txt}. Graph isomorphism is
handled by the BLISS backend exposed through igraph. Verification of
the six finite-domain certificates uses neither BLISS nor a SAT solver.
No GPU is required. Select the thread count to match the CPU resources
allocated by the cluster scheduler.

A full run writes $176\,590$ checkpoint files. Use storage with adequate
space and support for atomic file replacement. The total runtime depends
on the hardware, thread count, time limits and difficult instances; a
short pilot is not a reliable prediction of completion time. No
historical whole-cluster runtime is asserted in this supplement.

\section{Quick verification of the forty exceptional pairs}
\label{sec:quick}

\begin{minipage}{\linewidth}
Run:
\begingroup\footnotesize
\begin{verbatim}
python3 verify40.py
\end{verbatim}
\endgroup
\end{minipage}

On Windows, use \texttt{python} instead of \texttt{python3} when that is
the installed command. The program checks the fixed input identities,
the binding of case indices to the seed catalogue, the seed adjacency
data, the finite premises for the thirty-four counting exclusions, and
all six certificate trees. The required outcome is a successful process
exit and a report containing forty excluded cases, with zero unresolved
cases in this local set. The default report is
\file{results/verified40.json}; an alternative path may be selected by
\texttt{--report}.

For the six certificates the expected aggregate replay counts are
$7\,499$ search nodes, $198\,288$ explicit fan witnesses, and
$448\,770$ domain deletions. Each fan witness specifies a centre and
three or four disjoint pairs of peripheral vertices. The checker tests
the colours of all spokes and matching edges directly. It also checks
every domain deletion and every branch of the contradiction tree.
An unknown edge in a partial colouring belongs to neither colour.

The thirty-four counting exclusions and six tree certificates partition
the same fixed set of forty cases. A passed local check must not be
interpreted as a passed full-catalogue check. The full-run combination
procedure in Section~\ref{sec:combine} checks this coverage separately.

\section{Compiling and checking the search program}
\label{sec:build}

\subsection{CMake build}
From the archive root, run:
\begingroup\footnotesize
\begin{verbatim}
cmake -S src -B build -DCMAKE_BUILD_TYPE=Release
cmake --build build --config Release --parallel 4
\end{verbatim}
\endgroup
With a single-configuration Linux generator, the executable is
\file{build/fan_ramsey17}. Test it and inspect the fixed catalogue:
\begingroup\footnotesize
\begin{verbatim}
./build/fan_ramsey17 self-test
./build/fan_ramsey17 catalog --data data
\end{verbatim}
\endgroup
With the Visual Studio generator on Windows, use
\file{build\Release\fan_ramsey17.exe} in place of
\file{./build/fan_ramsey17}. Run the Windows commands in a developer
terminal with the x64 C++ toolchain available.

\subsection{Direct builds}
On Linux, the alternative command
\begingroup\footnotesize
\begin{verbatim}
bash src/build_linux.sh
\end{verbatim}
\endgroup
creates \file{src/fan_ramsey17}. In the remaining examples replace
\file{./build/fan_ramsey17} by \file{./src/fan_ramsey17} when using this
build.

\begin{minipage}{\linewidth}
In the x64 Native Tools Command Prompt for Visual Studio, the equivalent
Windows procedure is:
\begingroup\footnotesize
\begin{verbatim}
pushd src
call build_windows.bat
popd
src\fan_ramsey17.exe self-test
src\fan_ramsey17.exe catalog --data data
\end{verbatim}
\endgroup
\end{minipage}

Here the executable for subsequent commands is
\file{src\fan_ramsey17.exe}. Use ordinary single-line commands in
Windows; a backslash used for Linux line continuation is not a Windows
line-continuation character.

\subsection{Expected catalogue dimensions}

\begin{center}
\begin{tabular}{@{}lr@{}}
\toprule
Quantity & Expected value \\
\midrule
Eight-vertex seeds & $8\,812$ \\
Nine-vertex seeds & $115\,174$ \\
Eight-vertex profiles & $6\,860$ \\
Nine-vertex profiles & $60\,041$ \\
Retained ordered seed pairs & $88\,294\,674$ \\
Cases per full task & $500$ \\
Number of tasks & $176\,590$ \\
Cases in the last task & $174$ \\
\bottomrule
\end{tabular}
\end{center}

The executable rejects a catalogue or seed file that differs from its
fixed dataset. It also checks catalogue dimensions and basic profile
fields. These checks are useful consistency tests, not an independent
proof of exhaustive generation.

For a small executable test that does not start the full search, run:
\begingroup\footnotesize
\begin{verbatim}
./build/fan_ramsey17 run --data data \
  --machine 1 --machines 1 --threads 1 \
  --task-ids 0 --case-budget 5 --seconds 5 \
  --symmetry on --results results/smoke
\end{verbatim}
\endgroup
This deliberately processes at most five cases. Use a different results
directory for a new independent full computation.

\section{Running a full computation on one machine}

\begingroup\footnotesize
\begin{verbatim}
./build/fan_ramsey17 run --data data \
  --machine 1 --machines 1 --threads 4 \
  --symmetry on --seconds 5 --results results/machine01
\end{verbatim}
\endgroup

The \texttt{--seconds} option is a time limit for each case, not for
the whole job. The \texttt{--threads} option controls concurrent
workers on this machine. The symmetry option activates the symmetry
constraints described in the article. An independent run with
\texttt{--symmetry off} is supported but may be slower. The default
task order, \texttt{spread}, is a deterministic permutation of all
assigned tasks and is not a sampling rule.

\begin{minipage}{\linewidth}
The status fields have the following meanings:
\begin{center}
\begin{tabular}{@{}p{0.16\linewidth}p{0.77\linewidth}@{}}
\toprule
Field & Interpretation \\
\midrule
\texttt{PENDING} & Cases not yet processed. \\
\texttt{UNSAT} & Cases for which the SAT solver reports that no
permitted completion exists. \\
\texttt{UNKNOWN} & Cases for which the time-limited attempt did not
decide satisfiability. These are not exclusions. \\
\texttt{SAT} & Cases with a candidate completion; the program checks
and stores its red adjacency data. \\
\bottomrule
\end{tabular}
\end{center}

\end{minipage}

Each task progress line reports counts within that task. Its
\texttt{wall\_seconds} is elapsed time since the current run began;
it is not the duration of that task alone. The final run line reports
the number of attempts in the invocation. The runner may finish
normally with unresolved cases because every assigned case has received
its time-limited attempt.

The flag \texttt{proof\_verified=false} is intentional in production
solver output: the runner does not emit and independently check a
propositional proof for every exclusion. A solver's negative result,
a repeated negative result and a replayed certificate are different
forms of evidence.

\section{Resuming, retrying and merging}
\label{sec:resume}

Repeating a \texttt{run} command with the same results directory resumes
pending work. Already resolved cases are skipped. Previously unresolved
cases are attempted again only if \texttt{--retry-unknown} is present.
For example, after the first pass has stopped:
\begingroup\footnotesize
\begin{verbatim}
./build/fan_ramsey17 run --data data \
  --machine 1 --machines 1 --threads 4 \
  --symmetry on --seconds 60 --retry-unknown \
  --results results/machine01
\end{verbatim}
\endgroup
Further passes may use $600$ or $3600$ seconds per case. These are
example retry settings, not a guarantee that every case will finish.
Keep the dataset and seed ordering unchanged between passes.

After a run has stopped, create a summary and a merged checkpoint set:
\begingroup\footnotesize
\begin{verbatim}
./build/fan_ramsey17 merge --data data \
  --inputs results/machine01 --output results/merged_01
\end{verbatim}
\endgroup
The output directory must be new or contain no checkpoint files. The
merge writes \file{summary.json}, \file{pending_tasks.json},
\file{unknown_tasks.json}, and the merged task files. A missing task is
counted as pending; duplicate completed entries are reported. The merge
rejects inconsistent positive and negative decisions for a case.

The field \texttt{all\_pairs\_resolved} means that both the pending
and unresolved counts are zero. It alone does not mean all cases have
been excluded: the positive count must also be checked. The field
\texttt{solver\_reported\_exclusion} records whether the solver has
reported a negative decision for every case. It does not assert
independent proof verification.

To request an orderly stop, issue this command from another terminal:
\begingroup\footnotesize
\begin{verbatim}
./build/fan_ramsey17 stop --results results/machine01
\end{verbatim}
\endgroup
Allow the running process to finish and return to the shell before
merging or copying its checkpoints. The next explicit run clears the
stop request. Lock directories such as \file{RUN.lock} and the
individual task locks protect
against concurrent writers. If a process was killed, a lock may remain;
remove only the stale lock after confirming that its process has
terminated. Do not delete checkpoints to recover a stale lock, and do
not let two running processes write the same results directory.

\section{Distributed computation on six machines}

Use exactly the same code and data on every machine. Select
\texttt{--machines 6} everywhere and a distinct \texttt{--machine}
number from 1 through 6. Task $t$, numbered from zero, belongs to machine
$(t\bmod 6)+1$. Each machine writes to its own directory. For machine
1, run:
\begingroup\footnotesize
\begin{verbatim}
./build/fan_ramsey17 run --data data \
  --machine 1 --machines 6 --threads 4 \
  --symmetry on --seconds 5 --results results/machine01
\end{verbatim}
\endgroup
On machines 2 through 6, change both the machine number and the final
directory name accordingly, for example \texttt{--machine 2} and
\file{results/machine02}. The six workers need not start together.
The number of machines is independent of the number of threads per
machine. For scheduled cluster jobs, request the appropriate CPU
allocation and run the corresponding command inside the batch job.
Scheduler directives depend on the local cluster and are not hard-coded
into the distribution.

Retry on each machine using its own existing directory and machine
number, with \texttt{--retry-unknown} and a larger time limit. Once all
processes have stopped, collect the six directories on one machine,
preserving their distinct names, and merge:
\begingroup\footnotesize
\begin{verbatim}
./build/fan_ramsey17 merge --data data \
  --inputs results/machine01 results/machine02 \
           results/machine03 results/machine04 \
           results/machine05 results/machine06 \
  --output results/merged_six_01
\end{verbatim}
\endgroup
If a shared filesystem is used, the same six distinct directories can
be merged in place after the workers stop. Do not merge only the
individual summary files: case-level checkpoints are required.

To change the number of machines, first stop all writers and merge their
checkpoints. Work may then be redistributed by the new task-modulo rule.
Each new worker needs the existing checkpoint files for its assigned
tasks in its own results directory. Preserve the original merged copy
until the redistributed run has been checked. Reassigning tasks changes
the allocation of work, not the case indices or mathematical search.

\section{Combining a full run with the local certificates}
\label{sec:combine}

After creating a complete merged checkpoint directory, run:
\begingroup\footnotesize
\begin{verbatim}
python3 combine_results.py --results results/merged_01 \
  --output results/combined.json
\end{verbatim}
\endgroup
Use \file{results/merged_six_01} instead for the six-machine example.
This program checks the checkpoint dataset, task ranges, status integrity
and case-level coverage, and reruns the forty-case verifier. It requires
zero pending cases and zero positive cases. Every unresolved case in
the merged run must belong to the fixed set excluded by the local
certificates. An unresolved case outside that set causes failure and
requires further computation.

For the reported fourth-pass counts, the expected combination is
\[
88\,294\,634\ \text{solver exclusions}
+34\ \text{counting exclusions}
+6\ \text{certificate exclusions}
=88\,294\,674.
\]
The combined report records the actual split in the reader's run; it
need not reproduce the same intermediate unresolved counts. It does
not replace old unresolved status values with solver exclusions and
does not label the whole computation as independently certified.

The previously reported cumulative solver counts are listed below for
comparison. They are not checkpoint data and cannot substitute for
the case-level coverage check.
\begin{center}
\begin{tabular}{@{}lrr@{}}
\toprule
Pass & Excluded by the solver & Unresolved \\
\midrule
First & $88\,294\,434$ & $240$ \\
Second & $88\,294\,572$ & $102$ \\
Third & $88\,294\,620$ & $54$ \\
Fourth & $88\,294\,634$ & $40$ \\
\bottomrule
\end{tabular}
\end{center}
Each of these summaries reported zero pending cases and zero positive
cases. Different hardware and timing can change the intermediate
unresolved set without changing the underlying instances.

\section{Independent checks and witness verification}

The independent Python data validator uses the standard library:
\begingroup\footnotesize
\begin{verbatim}
python3 tools/search17_validate.py self-test
python3 tools/search17_validate.py data --data data \
  --profile-samples 4 --case-samples 12 \
  --output results/data_check.json
\end{verbatim}
\endgroup
It checks every seed record and the basic catalogue fields, with full
profile and indexing checks on the selected samples. The sample counts
can be increased. This command is not a second exhaustive canonical
generation and does not check all negative solver decisions.

If a run reports a positive case, preserve the stored witness and
verify it before drawing a conclusion:
\begingroup\footnotesize
\begin{verbatim}
python3 tools/search17_validate.py verify \
  --input results/machine01/SAT_FOUND.json --order 17 \
  --output results/witness_check.json
\end{verbatim}
\endgroup
The program directly checks the proposed colouring for forbidden fans.
A valid admissible colouring on seventeen vertices would contradict
the proposed exclusion, so a positive result must never be discarded
because it differs from the reported computation.

For a single seed pair, such as one of the six certified cases, use:
\begingroup\footnotesize
\begin{verbatim}
./build/fan_ramsey17 pair --data data \
  --left-index 3477 --right-index 1235 \
  --seconds 60 --symmetry off --output results/pair.json
\end{verbatim}
\endgroup
The pair command selects indices in the seed files, not a task number.
Its time-limited result is separate from replay of the supplied proof
certificate for that pair.

\section{Regenerating the seed libraries and catalogue}
\label{sec:regenerate}

Carry out regeneration in a working copy. Use an isolated Python
environment, for example on Linux:
\begingroup\footnotesize
\begin{verbatim}
python3 -m venv .venv
. .venv/bin/activate
python -m pip install -r tools/requirements.txt
mkdir -p regenerated
\end{verbatim}
\endgroup
The requirements fix igraph 1.0.0, NumPy 2.3.5, texttable 1.7.0 and
NetworkX 3.7. Do not use
Python's optimisation flag \texttt{-O}: some preparation consistency
checks are assertions.

\subsection{Canonical generation through order nine}
\begingroup\footnotesize
\begin{verbatim}
python tools/canonical_fans.py --validate
python tools/canonical_fans.py --max-order 9 --target 17 --audit \
  --export regenerated/fan9_target17.graph6
\end{verbatim}
\endgroup
The first command tests the unrestricted generator against known small
graph counts. The second runs the restricted canonical augmentation
used in this search. No finite time limit is imposed here; the program
refuses to export an incomplete level as a full seed library.
The numbers of retained isomorphism classes at orders $0$ through $9$
are expected to be
\[
1,\ 1,\ 2,\ 4,\ 11,\ 34,\ 155,\ 941,\ 8\,812,\ 115\,174.
\]
The restrictions are hereditary fan and clique exclusions. A minimum
red degree of seven is imposed only on a completed seventeen-vertex
colouring, not on a seed during augmentation.

\subsection{Profiles, filters and stable case indexing}
To regenerate the eight-vertex library and construct the profile and
pair catalogue using the newly generated nine-vertex library, run:
\begingroup\footnotesize
\begin{verbatim}
python tools/search17_prepare.py --data regenerated/data \
  --source-nine regenerated/fan9_target17.graph6 \
  --canonical-source tools/canonical_fans.py
\end{verbatim}
\endgroup
The source-nine option is checked against the fixed nine-vertex seed
identity. As a less expensive separate test of profile preparation,
one may instead supply \file{data/fan9_target17.graph6}; this reuses the
archived nine-vertex library and does not regenerate that library.

Check the regenerated fixed inputs and the bindings of the forty
exceptional cases:
\begingroup\footnotesize
\begin{verbatim}
python tools/check_data.py --data regenerated/data \
  --cases certificates/cases.json
python tools/search17_validate.py data --data regenerated/data \
  --profile-samples 6 --case-samples 24 \
  --output regenerated/data_check.json
./build/fan_ramsey17 catalog --data regenerated/data
\end{verbatim}
\endgroup

The expected cumulative numbers of ordered pairs are:
\begin{center}
\begin{tabular}{@{}lr@{}}
\toprule
Stage & Ordered pairs retained \\
\midrule
All eight-by-nine seed pairs & $1\,014\,913\,288$ \\
Maximum-degree compatibility & $682\,271\,903$ \\
Cross-edge count intervals & $682\,057\,873$ \\
Joint centre-type compatibility & $165\,925\,659$ \\
Saturated maximum-degree boundary & $88\,294\,674$ \\
\bottomrule
\end{tabular}
\end{center}

The preparation writes a new catalogue, identity information and a
corresponding C++ identity header in its output tree. It also generates
profile caches and run reports for that invocation. These generated
caches are not required in the published archive. Check that the
regenerated fixed files agree with the archived identities before
using existing checkpoints. If the ordering changes, even when the
graphs are isomorphic and the counts agree, old case indices no longer
identify the same instances. Do not bypass the executable's identity
checks or mix checkpoints from different catalogues.

\section{Replaying and regenerating certificates}

For the ordinary verification procedure, use Section~\ref{sec:quick};
certificate regeneration is optional. Each finite-domain certificate
uses the degree interval $7\leq d_R(v)\leq10$ and the necessary
one-for-one swap inequalities for a maximum red $8+9$ cut. Here
$d_R(v)$ is the red degree of vertex $v$. The checker validates these
finite constraints using the fixed seed adjacency data. The article
proves why every admissible seventeen-vertex colouring has such a
partition.

The six certificate case indices and their seed indices are:
\begin{center}
\begin{tabular}{@{}rrrr@{}}
\toprule
Case index & Eight-vertex seed & Nine-vertex seed & Search nodes \\
\midrule
$51\,996\,400$ & $3477$ & $1235$ & $911$ \\
$64\,095\,209$ & $4228$ & $1235$ & $917$ \\
$69\,103\,064$ & $4877$ & $68943$ & $1$ \\
$69\,257\,548$ & $4893$ & $68943$ & $1$ \\
$71\,867\,158$ & $5479$ & $1235$ & $5668$ \\
$82\,561\,944$ & $6899$ & $30887$ & $1$ \\
\bottomrule
\end{tabular}
\end{center}

The CMake build in Section~\ref{sec:build} also builds the certificate
generator \file{build/fan_tail_certify}. Regenerate all local proof data
in a new directory with:
\begingroup\footnotesize
\begin{verbatim}
python3 tools/regenerate_certificates.py \
  --binary build/fan_tail_certify \
  --output generated_certificates --seconds 120
\end{verbatim}
\endgroup
For a Visual Studio build, use
\file{build/Release/fan_tail_certify.exe} as the binary path.
The helper rebuilds the thirty-four-case counting certificate,
selects the six remaining cases from the fixed input list, calls the
C++ generator without symmetry restrictions, and independently replays
all regenerated certificates. The output directory must be empty.
The time limit is per case; if a certificate is not produced before
the limit, use a larger limit and a new output directory.

\begin{minipage}{\linewidth}
The regenerated files can also be checked in a separate invocation:
\begingroup\footnotesize
\begin{verbatim}
python3 verify40.py \
  --structural generated_certificates/theory_34_certificate.json \
  --csp-dir generated_certificates \
  --report results/regenerated_verified40.json
\end{verbatim}
\endgroup
\end{minipage}


Certificate generation must disable the optional twin-row symmetry
restriction, because the replay checker verifies a complete branching
tree without invoking a symmetry argument. A timeout or incomplete
tree is not an exclusion certificate. Store newly generated
certificates in a fresh directory and replay them before using them.
The helper enforces this symmetry setting. The archived certificates
are sufficient to reproduce all six independent exclusions without
rebuilding the generator.

\section{Submission and online archive}

The article and this supplement are intended to accompany the journal
submission. The separate reproducibility archive contains the final
source code, fixed inputs, independent checkers and certificate data;
it also includes this guide so that the archive is self-contained.
Research notebooks, obsolete executables, intermediate checkpoints,
local dependency installations and development logs are not part of
that distribution. Required third-party licences are retained.

The planned public repository is Zenodo, \url{https://zenodo.org/}.
The record has not yet been deposited and no public record DOI is
asserted here. Zenodo permits a DOI to be reserved in a draft upload;
the authors can then insert that persistent record link in the
article and in this supplement before releasing the files. The DOI
of the supplementary archive is distinct from the DOI eventually
assigned to the article. See the official instructions at
\url{https://help.zenodo.org/docs/deposit/describe-records/reserve-doi/}.
% Replace the preceding planned-deposit wording with the actual record
% DOI and publication status when the authors create the Zenodo record.

\end{document}
